\chapter*{Заключение}
\addcontentsline{toc}{chapter}{Заключение}
\label{ch:conclusion}

В ходе выполнения практической работы~\textnumero\,\cfgLabNumber{} была настроена
виртуальная машина \texttt{gateway} под управлением Ubuntu~20.04~LTS,
выполняющая роль шлюза локальной сети.

В рамках работы выполнены следующие задачи:
\begin{itemize}
  \item настроены интерфейсы \texttt{enp0s3}~(WAN) и \texttt{enp0s8}
    (LAN, \texttt{192.168.\cfgLabStudentN.1/24}) через Netplan;
  \item включён IP-форвардинг (\texttt{net.ipv4.ip\_forward=1})
    и созданы правила NAT \texttt{MASQUERADE} через \texttt{iptables};
  \item настроен и запущен DHCP-сервер \texttt{isc-dhcp-server}
    с адресным пулом
    \texttt{192.168.\cfgLabStudentN.10}--\texttt{192.168.\cfgLabStudentN.254};
  \item клиентская машина \texttt{Desktop1} автоматически получила
    IP-адрес, маршрут и DNS через DHCP.
\end{itemize}

Связь \texttt{Desktop1} с шлюзом и доступ в Интернет проверены командами
\texttt{ping~192.168.\cfgLabStudentN.1} и \texttt{ping~ya.ru}.
Полученные навыки~--- настройка маршрутизации и NAT, администрирование
DHCP-сервера~--- являются базовыми компетенциями сетевого администратора.
